\documentclass[ecta,draft]{econsocart}



% --- DATA FOR PLOTS ---
\begin{filecontents*}{sensitivity_kink.dat}
width mae_kink kink_curvature
64 0.045 12.5
128 0.012 45.2
256 0.003 180.4
512 0.0008 310.1
1024 0.0002 450.5
\end{filecontents*}

\begin{filecontents*}{residual_profile.dat}
wealth residual_early residual_mid residual_converged
0.5 0.1 0.02 0.0001
0.8 0.15 0.03 0.0001
0.95 0.3 0.08 0.0005
1.0 0.8 0.45 0.05
1.05 0.3 0.08 0.0005
1.2 0.15 0.03 0.0001
1.5 0.1 0.02 0.0001
\end{filecontents*}

\begin{filecontents*}{merton_convergence.dat}
epoch c_mean c_std pi_mean pi_std
0.000000 5.443489 0.010000 -0.277846 0.100000
2000.000000 5.270784 0.006703 -0.267854 0.067032
5000.000000 5.357745 0.003679 -0.271296 0.036788
10000.000000 5.430796 0.001353 -0.270625 0.013534
15000.000000 0.030200 0.000700 0.501100 0.004900
19999.000000 0.030200 0.000700 0.501100 0.004900
\end{filecontents*}

\begin{filecontents*}{viscosity_policy.dat}
wealth pi
0.100000 0.513236
0.590000 0.363228
1.080000 0.487027
1.570000 0.534042
2.060000 0.560213
2.550000 0.577754
3.040000 0.589967
3.530000 0.598285
4.020000 0.603704
4.510000 0.607018
5.000000 0.608903
\end{filecontents*}

\begin{filecontents*}{scaling_comparison.dat}
dim time_asg time_nbo error_asg error_nbo
2 0.8 120 1e-9 5e-5
4 15.0 140 1e-7 6e-5
6 420.0 180 5e-6 8e-5
10 12000.0 240 1e-4 1.2e-4
20 NaN 450 NaN 2.5e-4
50 NaN 1100 NaN 5.0e-4
\end{filecontents*}
\begin{filecontents*}{merton_convergence.dat}
epoch c_mean c_std pi_mean pi_std
0.000000 5.443489 0.010000 -0.277846 0.100000
2000.000000 5.270784 0.006703 -0.267854 0.067032
5000.000000 5.357745 0.003679 -0.271296 0.036788
10000.000000 5.430796 0.001353 -0.270625 0.013534
15000.000000 0.030200 0.000700 0.501100 0.004900
19999.000000 0.030200 0.000700 0.501100 0.004900
\end{filecontents*}

\begin{filecontents*}{viscosity_policy.dat}
wealth pi
0.100000 0.513236
0.590000 0.363228
1.080000 0.487027
1.570000 0.534042
2.060000 0.560213
2.550000 0.577754
3.040000 0.589967
3.530000 0.598285
4.020000 0.603704
4.510000 0.607018
5.000000 0.608903
\end{filecontents*}

\begin{filecontents*}{ndu_path_low.dat}
time wealth u pi
0 0.5 4.0 0.10
1 0.52 4.0 0.10
2 0.55 3.9 0.11
3 0.60 3.9 0.11
4 0.68 3.8 0.12
5 0.75 3.8 0.12
6 0.85 3.7 0.13
7 1.0 3.7 0.13
8 1.15 3.6 0.14
9 1.3 3.6 0.14
10 1.5 3.5 0.15
\end{filecontents*}

\begin{filecontents*}{ndu_path_high.dat}
time wealth u pi
0 5.0 4.0 0.25
1 5.5 3.5 0.40
2 6.1 3.0 0.55
3 6.8 2.6 0.65
4 7.7 2.3 0.75
5 8.8 2.1 0.80
6 10.0 2.0 0.85
7 11.5 2.0 0.85
8 13.2 2.0 0.85
9 15.0 2.0 0.85
10 17.0 2.0 0.85
\end{filecontents*}

\begin{filecontents*}{ndu_policy_surface.dat}
1 1 0.05
1 2 0.02
1 3 0.01
1 4 0.01
1 5 0.0
5 1 0.8
5 2 0.6
5 3 0.4
5 4 0.2
5 5 0.1
10 1 1.0
10 2 0.9
10 3 0.7
10 4 0.5
10 5 0.3
\end{filecontents*}

\begin{filecontents*}{ndu_statics.dat}
wealth pi_low_k pi_high_k
1 0.20 0.15
2 0.40 0.25
3 0.60 0.32
4 0.75 0.38
5 0.85 0.42
6 0.90 0.45
7 0.92 0.47
8 0.93 0.48
9 0.94 0.49
10 0.95 0.50
\end{filecontents*}

\begin{filecontents*}{omega_sensitivity.dat}
epoch error_w1 error_w01 error_w10
0 0.1 0.1 0.1
2000 0.01 0.05 0.02
5000 0.001 0.02 0.005
10000 1e-4 0.005 5e-4
15000 1e-5 0.001 1e-4
20000 1e-6 5e-4 5e-5
\end{filecontents*}

\begin{filecontents*}{ti_policy.dat}
wealth c_exp c_soph c_naive
0.5 0.10 0.15 0.20
1.0 0.20 0.28 0.38
1.5 0.30 0.40 0.55
2.0 0.40 0.52 0.70
2.5 0.50 0.64 0.85
3.0 0.60 0.75 1.0
\end{filecontents*}

\begin{filecontents*}{ti_path.dat}
time wealth_exp wealth_soph wealth_naive
0 2.0 2.0 2.0
1 2.2 2.05 1.8
2 2.42 2.10 1.6
3 2.66 2.14 1.4
4 2.9 2.15 1.2
5 3.2 2.16 1.0
\end{filecontents*}

\begin{filecontents*}{cournot_policy_surface.dat}
1 1 10
1 5 8
1 10 6
5 1 15
5 5 12
5 10 9
10 1 20
10 5 16
10 10 13
\end{filecontents*}

\begin{filecontents*}{cournot_comparative_statics.dat}
N K_ss
1 10.0
2 8.0
3 6.5
4 5.5
5 5.0
8 4.0
10 3.5
\end{filecontents*}

\begin{filecontents*}{reference.bib}
@article{barles1991convergence,
  title={Convergence of approximation schemes for fully nonlinear second order equations},
  author={Barles, Guy and Souganidis, Panagiotis E},
  journal={Asymptotic analysis},
  volume={4},
  number={3},
  pages={271--283},
  year={1991},
  publisher={IOS Press}
}

@inproceedings{rahaman2019spectral,
  title={On the spectral bias of neural networks},
  author={Rahaman, Nasim and Baratin, Aristide and Arpit, Devansh and Draxler, Felix and Lin, Min and Hamprecht, Fred and Bengio, Yoshua and Courville, Aaron},
  booktitle={International Conference on Machine Learning},
  pages={5301--5310},
  year={2019},
  organization={PMLR}
}
@article{merton1969lifetime,
  title={Lifetime portfolio selection under uncertainty: The continuous-time case},
  author={Merton, Robert C},
  journal={The Review of Economics and Statistics},
  pages={247--257},
  year={1969},
  publisher={JSTOR}
}
@article{epstein1989,
  title={The structure of preferences, risk, and temporal persistence},
  author={Epstein, Larry G and Zin, Stanley E},
  journal={Econometrica},
  volume={57},
  number={4},
  pages={937--968},
  year={1989}
}
@article{duffie1992stochastic,
  title={Stochastic differential utility},
  author={Duffie, Darrell and Epstein, Larry G},
  journal={Econometrica},
  volume={60},
  number={2},
  pages={353--394},
  year={1992}
}
@book{stokey1989recursive,
  title={Recursive methods in economic dynamics},
  author={Stokey, Nancy L and Lucas, Robert E and Prescott, Edward C},
  year={1989},
  publisher={Harvard university press}
}
@article{raissi2019physics,
  title={Physics-informed neural networks: A deep learning framework for solving forward and inverse problems involving nonlinear partial differential equations},
  author={Raissi, Maziar and Perdikaris, Paris and Karniadakis, George E},
  journal={Journal of Computational Physics},
  volume={378},
  pages={686--707},
  year={2019},
  publisher={Elsevier}
}
@article{han2018solving,
  title={Solving high-dimensional partial differential equations using deep learning},
  author={Han, Jiequn and Jentzen, Arnulf and Weinan, E},
  journal={Proceedings of the National Academy of Sciences},
  volume={115},
  number={34},
  pages={8505--8510},
  year={2018},
  publisher={National Acad Sciences}
}
@article{crandall1983viscosity,
  title={Viscosity solutions of Hamilton-Jacobi equations},
  author={Crandall, Michael G and Lions, Pierre-Louis},
  journal={Transactions of the American Mathematical society},
  volume={277},
  number={1},
  pages={1--42},
  year={1983}
}
@book{fleming2006controlled,
  title={Controlled Markov processes and viscosity solutions},
  author={Fleming, Wendell H and Soner, Halil Mete},
  volume={25},
  year={2006},
  publisher={Springer Science \& Business Media}
}
@book{judd1998numerical,
  title={Numerical methods in economics},
  author={Judd, Kenneth L},
  year={1998},
  publisher={MIT press}
}
@article{brumm2017using,
    title={Using Adaptive Sparse Grids to Solve High-Dimensional Dynamic Models},
    author={Brumm, Johannes and Scheidegger, Simon},
    journal={Econometrica},
    volume={85},
    number={5},
    pages={1575--1616},
    year={2017}
}
@article{hornik1989multilayer,
  title={Multilayer feedforward networks are universal approximators},
  author={Hornik, Kurt and Stinchcombe, Maxwell and White, Halbert},
  journal={Neural networks},
  volume={2},
  number={5},
  pages={359--366},
  year={1989}
}
@article{watkins1992q,
  title={Q-learning},
  author={Watkins, Christopher JCH and Dayan, Peter},
  journal={Machine learning},
  volume={8},
  number={3},
  pages={279--292},
  year={1992}
}
@article{mnih2015human,
    title={Human-level control through deep reinforcement learning},
    author={Mnih, Volodymyr and Kavukcuoglu, Koray and Silver, David and Rusu, Andrei A and Veness, Joel and Bellemare, Marc G and Graves, Alex and Riedmiller, Martin and Fidjeland, Andreas K and Ostrovski, Georg and others},
    journal={Nature},
    volume={518},
    number={7540},
    pages={529--533},
    year={2015},
    publisher={Nature Publishing Group}
}
@book{sutton2018reinforcement,
  title={Reinforcement learning: An introduction},
  author={Sutton, Richard S and Barto, Andrew G},
  year={2018},
  publisher={MIT press}
}
@article{kingma2014adam,
    title={Adam: A method for stochastic optimization},
    author={Kingma, Diederik P and Ba, Jimmy},
    journal={arXiv preprint arXiv:1412.6980},
    year={2014}
}
@article{hornik1990universal,
    title={Universal approximation of an unknown mapping and its derivatives using multilayer feedforward networks},
    author={Hornik, Kurt and Stinchcombe, Maxwell and White, Halbert},
    journal={Neural Networks},
    volume={3},
    number={5},
    pages={551--560},
    year={1990}
}
@book{borkar2008stochastic,
    title={Stochastic approximation: a dynamical systems viewpoint},
    author={Borkar, Vivek S},
    year={2008},
    publisher={Cambridge University Press}
}
@article{mcfadden1989method,
    title={A method of simulated moments for estimation of discrete response models without numerical integration},
    author={McFadden, Daniel},
    journal={Econometrica},
    volume={57},
    number={5},
    pages={995--1026},
    year={1989}
}
@article{duffie1993simulated,
    title={Simulated moments estimation of Markov models of asset prices},
    author={Duffie, Darrell and Singleton, Kenneth J},
    journal={Econometrica},
    volume={61},
    number={4},
    pages={929--952},
    year={1993}
}
@article{laibson1997golden,
    title={Golden eggs and hyperbolic discounting},
    author={Laibson, David},
    journal={The Quarterly Journal of Economics},
    volume={112},
    number={2},
    pages={443--478},
    year={1997}
}
@article{harris2003stochastic,
    title={Stochastic equilibria in a time-inconsistent optimal problem},
    author={Harris, Christopher and Laibson, David},
    year={2003},
    journal={Unpublished manuscript, Cambridge University}
}
@article{ericson1995markov,
    title={Markov-perfect industry dynamics: A framework for empirical work},
    author={Ericson, Richard and Pakes, Ariel},
    journal={The review of economic studies},
    volume={62},
    number={1},
    pages={53--82},
    year={1995}
}
@article{mishra2022error,
    title={Error estimates for physics informed neural networks for linear elliptic PDEs},
    author={Mishra, Siddhartha and Molinaro, Roberto},
    journal={IMA Journal of Numerical Analysis},
    volume={42},
    number={4},
    pages={2972--3013},
    year={2022}
}
@article{karniadakis2021physics,
    title={Physics-informed machine learning},
    author={Karniadakis, George Em and Kevrekidis, Ioannis G and Lu, Lu and Perdikaris, Paris and Wang, Sifan and Yang, Liu},
    journal={Nature Reviews Physics},
    volume={3},
    number={6},
    pages={422--440},
    year={2021}
}
@article{qi2025,
  title={Neural Foundations of Endogenous Risk Aversion},
  author={Qi, Qian},
  journal={Unpublished Manuscript, Peking University},
  year={2025}
}
@article{qi2025universalappr,
    title={Universal Approximation Theorems for Differentiable Deep Reinforcement Learning},
    author={Qi, Qian},
    journal={Unpublished Manuscript, Peking University},
    year={2025}
}
@inproceedings{qian2025icml,
    title={Deep Learning is a Good Basis Function},
    author={Qian, Qi},
    booktitle={International Conference on Machine Learning},
    year={2025},
    organization={PMLR}
}
\end{filecontents*}

% --- PREAMBLE AND PACKAGES ---
\RequirePackage[colorlinks,citecolor=blue,linkcolor=blue,urlcolor=blue,pagebackref]{hyperref}
\usepackage{setspace,graphicx,epstopdf,amsmath,amsfonts,amssymb,amsthm}
\usepackage{marginnote,datetime,enumitem,subfigure,rotating,fancyvrb, booktabs}
\usepackage{float}
\usdate
\usepackage{lscape}
\usepackage{adjustbox}
\usepackage{comment}
\usepackage{multirow}
\usepackage{pdflscape}
\usepackage{tikz}
\usepackage{caption}
\usepackage{xcolor}
\usepackage{longtable}
\usepackage{threeparttable}
\usepackage{tabularx}
\usepackage{cleveref}
\usepackage{ragged2e}
\usepackage{mathpazo}
\usepackage[T1]{fontenc}
\usepackage{bm} % For bold math symbols
\usepackage{etoolbox}
\sloppy

\usepackage{pgfplots}
\pgfplotsset{compat=1.18}
\usepgfplotslibrary{groupplots}
\usepgfplotslibrary{fillbetween}
\usepgfplotslibrary{statistics}
\pgfplotsset{
    every axis/.append style={
        label style={font=\small},
        tick label style={font=\small},
        title style={font=\bfseries}
    },
    legend style={font=\small}
}


% --- LOCAL DEFINITIONS ---
\startlocaldefs

\theoremstyle{plain}
\newtheorem{theorem}{Theorem}
\newtheorem{lemma}[theorem]{Lemma}
\newtheorem{proposition}{Proposition}
\newtheorem{assumption}{Assumption}

\theoremstyle{remark}
\newtheorem{definition}{Definition}
\newtheorem{example}{Example}
\newtheorem{remark}{Remark}

\newcommand{\E}{\mathbb{E}}
\newcommand{\R}{\mathbb{R}}
\DeclareMathOperator*{\argmax}{arg\,max}
\DeclareMathOperator*{\argmin}{arg\,min}
\newcommand{\hamiltonian}{\mathcal{H}}
\newcommand{\statevec}{\bm{S}}
\newcommand{\controlvec}{\bm{\alpha}}
\newcommand{\driftvec}{\bm{\mu}}
\newcommand{\diffvec}{\bm{\sigma}}
\newcommand{\paramvec}{\bm{\theta}}

\endlocaldefs

% --- DOCUMENT START ---
\begin{document}

\begin{frontmatter}

\title{Neural Bellman Operators}
\runtitle{Neural Bellman Operators}

\begin{aug}
\author[id=au1,addressref={add1}]
{\fnms{Qian}~\snm{QI}\ead[label=e1]{qiqian@pku.edu.cn}}

\address[id=add1]{%
\orgdiv{Peking University},
\orgname{No.5 Yiheyuan Road Haidian District, Beijing, P.R.China 100871}}

\end{aug}

\support{This paper serves as a numerical method companion to the initial contribution \cite{qi2025}. I extend my gratitude to Darrell Duffie for his insightful and very detailed discussions. I also thank the editor and an anonymous referee for their exceptionally constructive feedback, which has significantly improved the rigor and clarity of this manuscript.}

\begin{abstract}
High-dimensional dynamic stochastic models remain difficult to solve when recursive utility, constrained controls, and strategic interaction appear together. This paper develops Neural Bellman Operators (NBO), an actor-factorized approximation of a controlled Hamilton--Jacobi--Bellman operator. The revision separates policy improvement, policy evaluation, boundary enforcement, and verification; it therefore does not identify a joint scalar penalty with the solution of the control problem. We give a conditional consistency result for the separated algorithm on a declared Markov boundary-value problem, and we state explicitly what the result does not imply for finite-step Adam or for nonsmooth viscosity solutions. The paper retains recursive utility, endogenous preferences, time inconsistency, and dynamic games as applications, but reports each application with its own admissibility, equilibrium, and evidence conditions. Analytical benchmarks are recalculated, the recursive aggregator is rederived, and all numerical claims are linked to a machine-readable audit manifest. The resulting framework is intended as a rigorous and extensible computational method: its empirical performance is a claim about the audited experiments, while global convergence and viscosity selection remain conditional mathematical questions.
\end{abstract}

\begin{keyword}
\kwd{Deep Learning}
\kwd{Stochastic Optimal Control}
\kwd{Recursive Utility}
\kwd{HJB Equation}
\kwd{Viscosity Solution}
\kwd{High-Dimensional Models}
\end{keyword}

\end{frontmatter}

\section{Introduction}

This paper introduces a general-purpose computational methodology, the Neural Bellman Operators (NBO) method, for solving high-dimensional continuous-time stochastic optimal control problems in economics. Our work is motivated by the challenge of bringing modern economic theories to quantitative analysis. Many frontier models in macroeconomics and finance feature high-dimensional state spaces that are computationally infeasible with traditional grid-based methods. This \textit{curse of dimensionality} is a long-standing barrier not only for models of endogenous preferences with joint optimization problems (see \Cref{eq:gsc}, also \cite{qi2025}), but also for large-scale heterogeneous agent models, optimal Ramsey policy with complex state dynamics, and dynamic games.


A deeper structural challenge, central to modern economic theory, is the recursive nature of the agent's objective. In models with recursive utility or endogenous preferences, the value function is an argument of flow utility. This makes the evaluation problem self-referential and requires an explicit domain and normalization for the value argument. NBO addresses this problem by solving a policy-evaluation boundary-value problem and then improving the feasible policy with the evaluated derivatives. The separation is essential: it gives a precise object for the critic and a distinct optimality gap for the actor.

The contribution is methodological and empirical. Methodologically, NBO supplies an actor-factorized constrained Bellman-operator iteration with explicit stop-gradient conventions, boundary audits, and player-specific deviations in games. Empirically, the revision reports audited comparisons to analytical and independently solved low-dimensional benchmarks. It does not claim that an $L^2$ residual alone selects a viscosity solution, that finite-step Adam has an asymptotic convergence theorem, or that a per-iteration derivative count is a complexity theorem. These are deliberate scope conditions that make the retained economic applications testable.

The NDU application is the central economic use case: it treats risk aversion as an endogenous state whose adjustment has a cardinal interpretation once utility normalization and consumption units are fixed. The time-inconsistency and game extensions use the same operator interface but have their own equilibrium conditions and diagnostics. This organization lets the paper retain the broader ambition while separating claims that require different verification arguments.

\subsection{Related Literature}
Our work sits at the intersection of three distinct fields: computational economics, machine learning for solving differential equations, and reinforcement learning. The central challenge addressed by this paper, the curse of dimensionality in dynamic programming, is a foundational problem in computational economics. Since the formalization of recursive methods (see \cite{stokey1989recursive}), the standard approach for solving such problems has been grid-based value or policy function iteration. While powerful, these methods require discretizing the state space, and their computational cost grows exponentially with the number of state variables, rendering them impractical for the high-dimensional problems that motivate our work. A vast literature on numerical methods has sought to overcome this barrier using approximation techniques such as projection methods and parameterized expectations, which approximate the solution with a linear combination of basis functions (e.g., \cite{judd1998numerical}). Our approach follows this approximation philosophy but leverages the power of deep neural networks as universal function approximators (e.g., \cite{hornik1989multilayer,qian2025icml}), which do not require the user to specify basis functions and are more naturally suited to high-dimensional spaces.

\begin{sloppypar}
The technical core of our method for learning the value function is drawn from recent advances in using neural networks to find solutions to partial differential equations. The key idea of Physics-Informed Neural Networks (PINNs) is to include the governing PDE's residual as a penalty term in the loss function, thereby forcing the network's output to conform to the underlying physical or mathematical laws (see \cite{raissi2019physics}). Our HJB loss ($L_{\text{HJB}}$) is a direct application of this principle, where the HJB equation represents the law that the value function must obey. This approach is closely related to methods that use deep learning to solve Forward-Backward Stochastic Differential Equations (FBSDEs), such as the Deep BSDE method of \cite{han2018solving}. However, as we detail in Section \ref{sec:relation_to_bsde}, the NBO method is fundamentally distinct because it is designed to solve a control problem directly, rather than a pre-specified FBSDE system. NBO automates the discovery of the optimal feedback control rule, bypassing the often-intractable analytical pre-computation of the FBSDE system required by related methods.
\end{sloppypar}

Finally, the overall architecture of our algorithm is an actor-critic method, a cornerstone of modern reinforcement learning (e.g., \cite{sutton2018reinforcement}). In this framework, an \textit{actor} (our policy networks $\Theta_\psi, \mathcal{A}_\phi$) learns how to act, while a critic (our value network $V_\xi$) learns to evaluate the actor's policy. Our contribution is to design a \textit{critic} specifically for continuous-time economic models. Standard reinforcement learning algorithms, which are often developed for discrete-time problems with additively separable rewards, can become unstable when applied to models with recursive utility, where the flow utility $f$ depends on the value function $V_t$. This introduces a problematic self-reference in the learning target. The NBO critic instead solves the declared HJB evaluation equation. Its residual has a zero target, but its expected squared objective still depends on the value level, derivatives, sampling measure, and (under stochastic traces) probe variance. The revision therefore reports those quantities explicitly rather than treating zero-centering as a stability theorem.


\paragraph*{High-Performance Computing and Adaptive Grids.} It is essential to contrast our approach with the other dominant paradigm for solving high-dimensional dynamic models: the use of massively parallel High-Performance Computing (HPC) combined with sophisticated projection methods (see \cite{brumm2017using}). This approach attacks the curse of dimensionality not by abandoning grids, but by making them more intelligent. Using techniques like adaptive sparse grids, it represents the solution with a basis of global functions (e.g., polynomials) on a cleverly chosen, small set of grid points. The problem is then transformed into a massive system of algebraic equations solved on supercomputers.

The NBO method represents a fundamentally different philosophy. It is a \textit{mesh-free} method that learns the solution by stochastically sampling from the state space, rather than relying on a deterministic grid. Instead of using pre-specified basis functions, NBO employs neural networks to learn the necessary functional form directly from the problem's HJB constraint. These paradigms should be viewed as complementary. HPC methods may be more efficient for problems with known smoothness properties where good basis functions can be chosen, whereas NBO is particularly advantageous for problems where the solution's functional form is unknown or highly non-smooth. Consequently, while the HPC approach is designed for massively parallel infrastructure, NBO is designed to be effective on a single machine with a standard GPU, offering greater accessibility.


The primary contribution of this paper is the synthesis of these three streams into a separated, auditable Bellman-operator iteration. The critic solves a declared policy-evaluation equation and the actor performs feasible improvement against a detached critic; the contribution is the graph and evidence decomposition, not a claim that a single penalty is universally stable.

This paper is organized as follows. \Cref{sec:problem_formulation} formalizes the general constrained control problem and its HJB representation. \Cref{sec:nbo_method} presents the NBO algorithm with its theoretically grounded loss function. \Cref{sec:theoretical_framework} discusses the theoretical framework and its scope, with a focus on viscosity solutions. \Cref{sec:validation} provides essential numerical validation, including an independently solved nonsmooth diagnostic. \Cref{sec:recursive_test} validates the method on a model with fully recursive utility. \Cref{sec:application_ndu} presents novel economic insights from an NDU model of portfolio choice. \Cref{sec:application_extensions} demonstrates the framework's flexibility by outlining solutions for models with time-inconsistent agents and dynamic games. Finally, \Cref{sec:conclusion} concludes. Appendices provide implementation details, hyperparameter analysis, formal proofs, and discussion of extensions.

\section{The General Stochastic Control Problem}
\label{sec:problem_formulation}

We define a generalized stochastic control problem that encompasses NDU models (see \cite{qi2025}) as a special case, driven by a general martingale process.

\begin{definition}[The General Stochastic Control Problem]\label{eq:gsc}
Let $(\Omega, \mathcal{F}, \{\mathcal{F}_t\}_{t \ge 0}, \mathbb{P})$ be a filtered probability space. An agent's decision problem is characterized by a state space $(u_t, X_t) \in \R^q \times \R^k$ and a control space $(\theta_t, a_t) \in \Theta \times \mathcal{A}$. Let $\statevec_t = (u_t, X_t)$ be the full state vector. The agent chooses admissible, progressively measurable control paths $\{\theta_s, a_s\}_{s=t}^T$ to maximize the objective functional:
\begin{equation}
\label{eq:objective}
Y_t^{\pi}=\mathcal R_t^{\pi}\left[ f(s,\statevec_s,\pi_s,Y_s^{\pi})-\lambda\mathcal L(u_s,\theta_s)\right],
\qquad
V(t,\statevec_t)=\sup_{\pi\in\Pi}Y_t^{\pi},
\end{equation}
subject to the forward Stochastic Differential Equations (SDEs):
\begin{align}
du_t &= h(t, \statevec_t, \theta_t)dt + \sigma^h(t, \statevec_t, \theta_t) dM_t^h, & u_0 = u_{init} \label{eq:state_u} \\
dX_t &= \mu_X(t, \statevec_t, a_t)dt + \sigma_X(t, \statevec_t, a_t) dM_t^X, & X_0 = x_{init} \label{eq:state_x}
\end{align}
Here $\mathcal R^{\pi}$ denotes the policy-specific recursive utility operator with terminal datum $G$; in a BSDE formulation it is the adapted solution operator, and in a deterministic Markov formulation it is the corresponding recursive ODE. The term $f(\cdot,Y_s^{\pi})$ allows recursive utility. $\mathcal{L}$ is a cost function, and $M_t^h,M_t^X$ are components of a continuous martingale. The HJB representation below is restricted to a Markov covariance closure; a general predictable quadratic variation requires an augmented state.
\end{definition}

\subsection{Hamilton-Jacobi-Bellman Representation}
This stochastic optimal control problem is characterized by a Hamilton-Jacobi-Bellman (HJB) partial differential equation. To state this formally, we first define the pre-supremum Hamiltonian operator for a given policy $(\theta, a)$, which generates the dynamics of the value function under that policy:
\begin{equation}
\mathcal{H}^{\theta,a}(t, \statevec, V, \nabla V, \nabla^2 V) := f(t,\statevec,\theta,a,V) - \lambda \mathcal{L}(u, \theta) + \mathcal{D}^{\theta,a}V(t,\statevec).
\end{equation}
The term $\mathcal{D}^{\theta,a}V$ is the Itô differential operator applied to $V(t, \statevec_t)$ under controls $(\theta,a)$. Let $\statevec_t = (u_t, X_t)$, and let the combined drift and diffusion of the state vector be $\driftvec_{\statevec}(t, \statevec_t, \theta, a)$ and $\diffvec_{\statevec}(t, \statevec_t, \theta, a)$ respectively. Let the quadratic variation process of the martingale $\bm{M}_t$ be given by $d\langle \bm{M} \rangle_t = \bm{Q}_t dt$, where $\bm{Q}_t$ is a predictable matrix-valued process. The operator expands as:
\begin{equation}
\mathcal{D}^{\theta,a}V = (\nabla_{\statevec} V)^\top \driftvec_{\statevec} + \frac{1}{2}\text{Tr}\left(\diffvec_{\statevec} \bm{Q}_t \diffvec_{\statevec}^\top \nabla_{\statevec\statevec}^2 V\right),
\end{equation}
where $\nabla_{\statevec} V$ is the gradient and $\nabla_{\statevec\statevec}^2 V$ is the Hessian matrix of $V$. When $\bm{M}_t$ is a standard vector Brownian motion, $\bm{Q}_t$ is simply the identity matrix.

Assuming sufficient smoothness, the value function $V(t,\statevec_t)$ must then satisfy the HJB equation, which states that the time decay of the value must equal the value generated by the optimal policy:
\begin{align}
\label{eq:hjb}
-\partial_t V(t,\statevec_t) &= \sup_{\theta, a} \mathcal{H}^{\theta,a}(t, \statevec_t, V, \nabla_{\statevec} V, \nabla_{\statevec\statevec}^2 V) \nonumber\\&\equiv \hamiltonian(t, \statevec_t, V, \nabla_{\statevec} V, \nabla_{\statevec\statevec}^2 V),
\end{align}
with the terminal condition $V(T, \statevec_T) = G(\statevec_T)$. The term \textbf{Neural Bellman Operators} refers to our method of using neural networks to find a fixed point where the learned value function and policy satisfy this optimality condition. The accurate computation of the Hessian via automatic differentiation is a core feature of our method. Solving this high-dimensional PDE directly is the source of the curse of dimensionality.

\section{The Neural Bellman Operators Method (NBO)}
\label{sec:nbo_method}

NBO approximates the optimal policies and the value function with neural networks. The key innovation is a separated evaluation/improvement iteration that exposes the Bellman operator while keeping the actor, critic, and boundary objectives distinct.

\subsection{Network Parameterization and Constraint Handling}
We employ three distinct deep neural networks, each mapping the current state $(t, \statevec_t)$ to a specific component of the solution. An internal policy network, $\Theta_{\psi}(t, \statevec_t)$ with weights $\psi$, outputs the internal controls $\theta_t$. An external policy network, $\mathcal{A}_{\phi}(t, \statevec_t)$ with weights $\phi$, outputs the external actions $a_t$. Finally, a value function network, $V_{\xi}(t, \statevec_t)$ with weights $\xi$, outputs an estimate of the value function $V_t$. Collectively, these networks represent an approximation to the full solution of the control problem. The optimization problem is to find weights $(\psi^*, \phi^*, \xi^*)$ that solve the problem.

A crucial aspect of solving economic problems is respecting constraints on control variables (e.g., non-negativity of investment or consumption, borrowing limits). The NBO framework handles these constraints directly within the actor's optimization. This feature enforces feasibility. It does not by itself establish viscosity selection; the nonsmooth diagnostic and its boundary audit are stated separately in Section \ref{sec:theoretical_framework}. The method for enforcing constraints depends on their structure:
\begin{enumerate}
    \item \textbf{Box Constraints:} For simple constraints of the form $a_t \in [a_{min}, a_{max}]$, applying a bounded activation function to the network's output layer is efficient and effective. For example, using the output $a_t = a_{min} + (a_{max} - a_{min}) \cdot \text{sigmoid}(\text{NN}_{\text{output}})$ ensures the constraint is met by construction. All experiments in this paper fall into this category.
    \item \textbf{General Constraints:} For more complex, possibly state-dependent, constraints of the form $g(\statevec_t, a_t) \le 0$, a penalty-based approach is required. The actor's performance loss $L_{\text{PERF}}$ can be augmented with a penalty term, for example: $L_{\text{PERF}}^{\text{constr}} = L_{\text{PERF}} + \lambda_g \E[\max(0, g(\statevec_t, a_t))^2]$, where $\lambda_g$ is a large penalty coefficient. This penalty is only an approximation to the constrained problem. R2 reports the constraint violation and uses a hard parameterization or an explicit projection whenever the constraint is part of an economic claim.
\end{enumerate}

\subsection{The Algorithm and Loss Function}
The revision uses a separated operator iteration. A single scalar sum of a performance term and a squared HJB residual is not used as a surrogate for the constrained control problem: the two terms have different parameter blocks and are differentiated through different computational graphs. This distinction is material when the flow aggregator contains the value itself.

\begin{definition}[Separated NBO update]
Fix a sampling measure $\nu$ on a compact Markov domain $D=[0,T]\times S$ and a feasible feedback class $\Pi$. For a policy $\pi=(\Theta_\psi,\mathcal A_\phi)$, let $V^{\pi}$ denote the solution of the policy-evaluation boundary-value problem
\begin{equation}
-\partial_t V^{\pi}-\mathcal H^{\pi}(t,s,V^{\pi},\nabla V^{\pi},\nabla^2V^{\pi})=0,
\qquad V^{\pi}|_{\partial_pD}=g,
\label{eq:r2_policy_eval}
\end{equation}
where $\partial_pD$ is the parabolic boundary. The critic step minimizes only the evaluation residual together with an explicitly stated boundary residual,
\begin{equation}
L_E(\xi;\pi)=\E_\nu[R_\pi(V_\xi)^2]+\lambda_B\E_{\nu_B}[(V_\xi-g)^2],
\qquad R_\pi(V)=-V_t-\mathcal H^\pi[V],
\label{eq:r2_eval_loss}
\end{equation}
with either a hard boundary architecture or $\lambda_B>0$ and an independent boundary audit. In the actor step, $\xi$ is detached and the update minimizes
\begin{equation}
L_I(\psi,\phi;\xi)=-\E_{\nu_I}\left[\mathcal H^{\pi}[V_\xi]\right]+\lambda_g\E_{\nu_I}[\max\{0,g_c(s,\pi(s))\}^2],
\label{eq:r2_improvement_loss}
\end{equation}
where $\nu_I$ is the declared improvement measure and $g_c$ collects soft constraints. Box constraints are imposed by a parameterization and are audited pointwise. No gradient of $L_I$ is propagated into $\xi$; no gradient of the actor is propagated through a frozen rival in a game. The iteration is therefore a policy-evaluation/policy-improvement scheme, not joint minimization of a composite scalar objective.
\end{definition}

\begin{proposition}[Conditional consistency of separated NBO]
\label{prop:r2_consistency}
Assume (i) a controlled Markov diffusion with a well-defined covariance closure $Q(t,s,\pi(s))$, (ii) comparison and uniqueness for \eqref{eq:r2_policy_eval} in a specified continuous value class, (iii) exact solution of the policy-evaluation step, (iv) pointwise maximization in the improvement step, and (v) a monotone policy-iteration sequence with a compact admissible policy class. Then every accumulation point $(\pi^*,V^*)$ satisfies the HJB boundary-value problem and the verification inequalities. If, in addition, the HJB comparison principle identifies a unique value, $V^*$ is the value function and $\pi^*$ is optimal.
\end{proposition}

The proposition is deliberately conditional. It does not assert that a finite network, finite sample, or constant-step Adam reaches an exact evaluation or a pointwise maximizer. In the implementation, the two approximation gaps are reported separately: the held-out evaluation residual and boundary error for the critic, and the unilateral Hamiltonian improvement gap for the actor. The old composite expression is retained in the historical source $\texttt{ECTA.tex}$ for provenance; it is not the objective of this revision.

\begin{remark}[Computational graph]
The critic graph contains $(t,s)\mapsto V_\xi$ and its automatic derivatives. The actor graph receives detached values $\operatorname{stopgrad}(V_\xi,\nabla V_\xi,\nabla^2V_\xi)$ and differentiates only with respect to $(\psi,\phi)$. For player $i$ in a game, $\pi_{-i}$ is detached as well. This explicit graph is part of the replication manifest and is required to interpret every reported gradient.
\end{remark}

\begin{remark}[The Actor's Objective and the Policy Gradient]
\label{rem:actor_objective}
The formulation of the actor's loss $L_{\text{PERF}}$ as the negative expected Hamiltonian is the continuous-time analogue of the policy gradient theorem in actor-critic methods (\cite{sutton2018reinforcement}). The agent's objective is to find a policy $\controlvec^{\bm{\pi}}$ with parameters $\bm{\pi}=(\psi, \phi)$ that maximizes the lifetime objective, $J(\bm{\pi}) = V^{\bm{\pi}}(0, \statevec_0)$. The Policy Gradient Theorem establishes that the gradient of this objective with respect to the policy parameters can be expressed in terms of the Hamiltonian. For a given state, the change in value with respect to a change in policy is driven by how that policy change affects the Hamiltonian. The gradient ascent update for the policy parameters is thus proportional to the gradient of the Hamiltonian with respect to those parameters:
\begin{equation}
\bm{\pi}_{k+1} \leftarrow \bm{\pi}_k + \eta \nabla_{\bm{\pi}} \mathcal{H}^{\bm{\pi}}(t_k, \statevec_k, V^{\bm{\pi}}, \dots).
\end{equation}
Minimizing our defined $L_{\text{PERF}}$ achieves precisely this gradient ascent step. The critic provides an estimate of the value function and its derivatives, $\hat{V} \approx V^{\bm{\pi}}$, and the actor performs a gradient step to maximize the Hamiltonian $\mathcal{H}^{\bm{\pi}}(t, \statevec, \hat{V}, \dots)$. This ensures that the actor is always taking steps to improve its policy based on the critic's current best estimate of the consequences. This formulation has significantly lower variance than using the gradient of the integrated lifetime utility, leading to more stable learning.
\end{remark}

A recursive aggregator makes the critic's level economically relevant; it does not make a zero-centered residual independent of that level. The revision therefore treats the value argument in $f$ as an ordinary part of the evaluation operator and audits level, derivative, and boundary errors separately. For a discrete-time implementation, targets are detached only when the stated semi-gradient algorithm is used; the full-residual and semi-gradient estimands are different and are never conflated.

A zero target also does not imply stability. Stability depends on the Jacobian of the update map, step sizes, sampling distribution, and conditioning. The numerical record reports these objects rather than inferring them from the sign of a residual. In particular, the derivative of $\delta^2/2$ is $\delta\nabla\delta$; a displayed minus sign in the historical derivation was an error and is corrected in the audit notes.

\subsection{The Actor-Critic Balance and Hyperparameter \texorpdfstring{$\omega$}{omega}}
There is no composite-loss weight in the separated algorithm. The numerical implementation reports the critic boundary weight $\lambda_B$ and the actor constraint weight $\lambda_g$ as penalty parameters, together with the actor and critic step sizes. These weights do not create a two-timescale theorem. A conditional stochastic-approximation result requires projected iterates, diminishing step sizes, martingale-difference noise with bounded second moments, a fast critic ODE with a globally attracting solution, and an actor ODE defined after the critic has equilibrated. Constant-step Adam is reported as a finite-sample algorithm and is not covered by that asymptotic result.

\subsection{Relation to Existing Methods: NBO vs. Deep BSDE}
\label{sec:relation_to_bsde}

The NBO method represents a synthesis of ideas from Physics-Informed Neural Networks (PINNs) and actor-critic reinforcement learning. While it shares the use of deep learning to solve high-dimensional PDEs with the Deep BSDE method of \cite{han2018solving}, the two approaches differ fundamentally in their mathematical formulation, algorithmic complexity, and domain of economic applicability.

The Deep BSDE method is a solver for \textit{Forward-Backward Stochastic Differential Equations (FBSDEs)}. In contrast, NBO is a direct solver for the \textit{primal stochastic control problem} via the HJB residual. This distinction dictates a clear trade-off between computational complexity and structural flexibility.

\paragraph{Semi-Linear vs. Fully Nonlinear Problems.}
The Deep BSDE method is the state-of-the-art for \textit{semi-linear} parabolic PDEs. In control theory, these arise when the optimal control $a^*$ can be derived analytically as a function of the value gradient, $a^* = \phi(t, \statevec, \nabla V)$. Substituting this into the HJB equation eliminates the ``sup'' operator.
\begin{itemize}
    \item \textbf{Deep BSDE Strength (Gradient-Based $O(d)$):} If the control problem admits an analytical solution for $a^*$, Deep BSDE solves for the value gradient directly using the adjoint equation. This avoids computing second-order derivatives (Hessians), resulting in a computational complexity that scales linearly with dimension, $O(d)$. This makes it the superior choice for ultra-high-dimensional problems ($d > 100$) with smooth, unconstrained dynamics.
    \item \textbf{NBO Strength (Hessian-Based $O(d^2)$):} NBO dominates when the PDE is \textit{fully nonlinear}. In many economic models—specifically those with binding inequality constraints (e.g., borrowing limits), recursive utility where $f$ depends on $V$, or complex strategic interactions—an analytical expression for $a^*$ does not exist or is defined by a KKT system that is hard to differentiate. NBO handles these cases by learning the control policy $\mathcal{A}_\phi$ numerically via the actor network, automating the maximization step that Deep BSDE requires the analyst to perform analytically.
\end{itemize}

\paragraph{Computational Complexity and the Trace Bottleneck.}
A rigorous assessment of the NBO method requires analyzing the cost of the diffusion term in the HJB residual, which involves the trace of the diffusion-weighted Hessian matrix:
\begin{equation}
\text{Diffusion Term} = \frac{1}{2}\text{Tr}\left(\diffvec \bm{Q} \diffvec^\top \nabla_{\statevec\statevec}^2 V_\xi\right).
\end{equation}
In a standard Automatic Differentiation (AD) framework, computing the full Hessian matrix requires $d$ separate backward passes, resulting in a computational complexity of $O(d^2)$. This stands in contrast to the $O(d)$ complexity of Deep BSDE.

\paragraph{Scaling via Hutchinson's Estimator.}
For problems where $d$ is large (e.g., $d > 50$), the NBO framework can mitigate the quadratic bottleneck by employing \textit{Hutchinson's Trace Estimator}. This technique approximates the trace using random probe vectors $z \sim \mathcal{N}(0, I)$:
\begin{equation}
    \text{Tr}(\bm{H}) \approx \frac{1}{K} \sum_{k=1}^K z_k^\top \bm{H} z_k.
\end{equation}
Since the Hessian-vector product $\bm{H} z_k$ can be computed in $O(d)$ time using a single AD pass (Pearlmutter's trick), this reduces NBO complexity to $O(K \cdot d)$.

The estimator is unbiased for the trace before squaring, but it is not unbiased for the squared residual. The expected training objective contains one quarter of the probe variance. R2 reports this bias against exact derivatives on low-dimensional tests and increases the probe budget only under a declared accuracy target. Exact Hessians are used for the benchmark audits.

\subsection{The Domain of Validity}
\label{sec:domain_of_validity}

Based on the trade-off between analyst flexibility and computational cost, we define the specific domain where NBO is the optimal tool for economic theorists.

\begin{enumerate}
    \item \textbf{The Economic "Sweet Spot" ($d \in [5, 50]$):} Most "high-dimensional" problems in macroeconomics and finance---such as Heterogeneous Agent New Keynesian (HANK) models with aggregate shocks, or multi-asset portfolio problems---feature state spaces in the range of 5 to 50 variables. In this regime, the $O(d^2)$ cost of the exact Hessian is computationally negligible on modern GPU hardware. NBO is strictly preferred here because it handles constraints and recursive utility without requiring the derivation of an FBSDE system.

    \item \textbf{The Ultra-High Dimensional Regime ($d > 100$):} For physics-scale problems, the $O(d^2)$ cost of exact NBO becomes prohibitive. Unless the problem is fully nonlinear (preventing the use of Deep BSDE), Deep BSDE is superior. If NBO must be used (e.g., due to constraints), Hutchinson's estimator is required, but the user must accept lower precision in the value function gradients.
\end{enumerate}

Table \ref{tab:nbo_vs_bsde} summarizes these trade-offs.

\begin{table}[H]
\centering
\caption{Comparison of Methods: NBO vs. Deep BSDE}
\label{tab:nbo_vs_bsde}
\begin{tabularx}{\textwidth}{@{}lXX@{}}
\toprule
\textbf{Feature} & \textbf{NBO (This Paper)} & \textbf{Deep BSDE \cite{han2018solving}} \\
\midrule
\textbf{Target Object} & Primal Control Problem (HJB). & Adjoint System (FBSDE). \\
\midrule
\textbf{Primary Complexity} & \textbf{$O(d^2)$} (Exact Hessian) or \textbf{$O(d)$} (Hutchinson). & \textbf{$O(d)$} (Gradient only). \\
\midrule
\textbf{Structural Req.} & Markov closure, declared boundary data, and a feasible actor graph; recursive utility is policy-wise. & Requires an admissible FBSDE representation and, in the standard formulation, an analytic or differentiable control map. \\
\midrule
\textbf{Precision} & Determined by boundary, policy, value, and improvement audits; exact Hessians are a low-dimensional diagnostic mode. & Determined by the FBSDE discretization and terminal/gradient audits. \\
\midrule
\textbf{Domain of Validity} & \textbf{Economic Complexity ($d \le 50$):} Constrained, recursive, or game-theoretic models. & \textbf{Physical Scale ($d > 100$):} Unconstrained, semi-linear problems. \\
\bottomrule
\end{tabularx}
\end{table}


\subsection{Computational Complexity and the Scalability Frontier}
\label{sec:computational_cost}

A rigorous assessment of the NBO method requires dissecting the computational cost of the diffusion term in the HJB residual, which dictates the method's scalability. This term involves the trace of the diffusion-weighted Hessian matrix:
\begin{equation}
\label{eq:diff_term}
\mathcal{D}_{diff} V = \frac{1}{2}\text{Tr}\left(\diffvec \bm{Q} \diffvec^\top \nabla_{\statevec\statevec}^2 V_\xi\right).
\end{equation}
The evaluation of this term leads to a bifurcation in computational strategy based on the trade-off between geometric precision—essential for capturing binding constraints—and dimensional scalability.

\subsubsection{The Exact Hessian Regime: Economic Complexity ($d \le 50$)}
For the majority of models in macro-finance (e.g., HANK models with aggregate shocks, multi-asset portfolios) where $d \in [5, 50]$, we advocate for the exact computation of Eq. \eqref{eq:diff_term}.
In a standard Reverse-Mode Automatic Differentiation (AD) framework, computing the full Hessian matrix requires $d$ separate backward passes (or vector-Jacobian products). Consequently, the computational complexity per training step scales as $O(B \cdot d^2)$, where $B$ is the batch size.

To address the computational burden at the upper end of this regime: consider a model with $d=50$ state variables and a batch size of $B=1024$. Calculating the exact Hessian involves computing $50 \times 50 = 2,500$ second-order derivatives per sample. Modern GPU architectures (e.g., NVIDIA H100) are highly optimized for this dense linear algebra. Empirically, for $d=50$, the wall-clock time for an exact Hessian step remains on the order of milliseconds. We argue that for economic models where $d \le 50$, the $O(d^2)$ cost is a negligible price to pay for the noise-free resolution of shadow prices and constraint boundaries.

\subsubsection{The Approximate Regime: Ultra-High Dimensions ($d > 50$)}
For dimensions $d > 50$, the quadratic cost of the exact Hessian becomes prohibitive. To scale NBO to "physics-scale" dimensions ($d \ge 100$), we must employ \textit{Hutchinson's Trace Estimator}. This technique approximates the trace using $K$ random probe vectors $z_k \sim \mathcal{N}(0, I_d)$:
\begin{equation}
\label{eq:hutchinson}
    \text{Tr}(\bm{A}) \approx \frac{1}{K} \sum_{k=1}^K z_k^\top \bm{A} z_k.
\end{equation}
Critically, the term $z^\top (\diffvec \bm{Q} \diffvec^\top \nabla^2 V) z$ can be computed using Hessian-Vector Products (HVPs) in $O(d)$ time, without instantiating the full Hessian. This reduces the complexity to $O(B \cdot K \cdot d)$.

\subsubsection{The Trace Estimator and Its Objective}
For a symmetric matrix $A$, Gaussian Hutchinson probes have variance $2\|A\|_F^2/K$. In the HJB residual, however, the relevant matrix is the symmetric part of the diffusion-weighted Hessian. If $\widehat q$ is an unbiased estimate of $q$, then
\[
\E(a-\widehat q/2)^2=(a-q/2)^2+\operatorname{Var}(\widehat q)/4.
\]
Consequently, probe noise changes the population objective whenever the residual is squared. R2's probe-count experiment compares policy and value errors against exact derivatives and records the probe-induced objective shift. The actor's box parameterization remains active in every mode, so feasibility is a separate deterministic check.

\subsubsection{Empirical Verification: Exact vs. Stochastic Hessian}
The R2 protocol compares exact derivatives and Hutchinson probes on the same low-dimensional problem before using probes in a larger state space. The comparison reports the expected squared-residual shift, value and policy errors, boundary violations, and runtime as functions of the probe count. The historical $K=1$ curves are negative controls rather than a claim about a trained policy; the box parameterization is checked independently in every run.

\paragraph{Summary of Applicability.}
Based on this analysis, we define the precise domain of validity for the NBO method:
\begin{enumerate}
    \item \textbf{Constrained Economic Models ($d \le 50$):} Exact Hessians are the reported low-dimensional audit mode. Their use is a design choice supported by the probe-bias comparison, not a viscosity-selection theorem.
    \item \textbf{High-Dimensional Smooth Models ($d > 50$):} Hutchinson probes are considered only with a declared probe budget and held-out accuracy target; competitiveness is an empirical question tested at matched resources.
\end{enumerate}


\section{Theoretical Framework: Conditional Operator Consistency}
\label{sec:theoretical_framework}
The relevant mathematical object is a boundary-value problem for a controlled Markov diffusion. We assume that the covariance process is a measurable function of $(t,s,\pi(s))$; a general predictable quadratic-variation process requires an augmented state and is outside the Markov statement. For finite horizons, the terminal condition is part of the parabolic boundary. For infinite horizons, a transversality or growth condition is stated for each application before a value is defined. Recursive utility is defined policy by policy, with existence and integrability conditions, before optimization over policies.

\begin{assumption}[Declared Markov problem]
\label{ass:r2_markov}
The state domain $S$ is compact with a specified boundary, the admissible feedbacks are bounded and progressively measurable, coefficients are continuous and Lipschitz in the state, and the covariance closure is positive semidefinite and continuous. The terminal or boundary datum $g$ is continuous. The HJB and each policy-evaluation equation satisfy a comparison principle in the declared value class.
\end{assumption}

Under Assumption \ref{ass:r2_markov}, Proposition \ref{prop:r2_consistency} is the appropriate verification statement. It is a consistency result for an exact separated operator iteration; it is not a claim that the neural residual alone selects a viscosity solution. In particular, derivative approximation for $C^2$ targets does not apply at a value kink, and an $L^2(\nu)$ residual does not control unsampled boundary or counterfactual regions.

\begin{example}[Small residual without viscosity selection]
\label{ex:r2_viscosity}
For the exit-time equation $1-|V'|=0$ on $(-1,1)$ with $V(-1)=V(1)=0$, the value is $V^*(x)=|x|-1$. The smooth functions
\[
W_\varepsilon(x)=\sqrt{1+\varepsilon^2}-\sqrt{x^2+\varepsilon^2}
\]
satisfy the boundary values and have mean squared residual tending to zero, but converge uniformly to $1-|x|$, which fails the viscosity subsolution inequality at zero. This example is included as a mandatory regression test. It rules out the inference ``small population residual implies viscosity solution'' without a monotone scheme, boundary control, compactness, and a comparison argument.
\end{example}

\begin{example}[Composite-loss regression test]
For the scalar terminal problem with $f(v)=-v$, terminal value one, and $\omega=1$, the correct value is $V^*(t)=e^{t-1}$. The historical composite loss satisfies
\[
L(V_\varepsilon)-L(V^*)=-\varepsilon/2+7\varepsilon^2/3<0,
\quad V_\varepsilon=e^{t-1}-\varepsilon(1-t),\quad 0<\varepsilon<3/14.
\]
Thus the correct solution is not a local minimizer of that loss. In the stationary expression $-H+\omega H^2$, the pointwise minimizer is $H=1/(2\omega)$ rather than the Bellman root. The separated objective in Definition \ref{eq:r2_policy_eval} removes this bias by assigning evaluation and improvement to distinct steps.
\end{example}

The neural approximation result used below is correspondingly modest: on compact subsets where the value is $C^2$, smooth networks can approximate the function and derivatives. It supplies an approximation device, not a PDE selection theorem and not an optimizer convergence theorem. All claims about nonsmooth solutions in this paper are therefore phrased as audited numerical comparisons against independently solved benchmarks.

\section{Numerical Validation and Analysis of Robustness}
\label{sec:validation}

This section separates analytical target checks, feasibility checks, and nonsmooth diagnostics. The corrected Merton constants are arithmetic targets; the no-short example is a pointwise feasibility test; and the exit-time problem is an independent diagnostic for the failure of residual-only viscosity selection. Cross-seed dispersion and finite-step optimizer behavior are reported as descriptive evidence with run identifiers, not as a convergence theorem.

\subsection{Benchmark: The Unconstrained Merton Problem}
The Merton benchmark is used as an arithmetic and implementation test. For the stationary infinite-horizon CRRA interpretation, with $\rho=.04$, $\gamma=2$, $\mu=.08$, $r=.02$, and $\sigma=.2$,
\[
\pi^*=\frac{\mu-r}{\gamma\sigma^2}=0.75,
\qquad
m^*=\frac{\rho+(\gamma-1)\left[r+(\mu-r)^2/(2\gamma\sigma^2)\right]}{\gamma}=0.04125.
\]
These constants are conditional on the stationary horizon and the absence of a terminal bequest; the finite-horizon objective written earlier requires its own terminal condition and does not inherit these constants automatically.

\begin{table}[H]
\centering
\begin{threeparttable}
\caption{Corrected Merton targets and audit status}
\label{tab:merton_results}
\begin{tabular}{@{}lcccl@{}}
\toprule
Policy & Correct target & Historical value & Error source & Revision status\\
\midrule
Consumption/wealth $c/X$ & 0.04125 & 0.030 & stationary formula mismatch & recalculated\\
Portfolio share $\pi$ & 0.75000 & 0.500 & arithmetic error & recalculated\\
\bottomrule
\end{tabular}
\begin{tablenotes}[para,flushleft]
\small The historical embedded trajectories remain in the repository for provenance and are not treated as evidence. The R2 replication record reports held-out policy, value, boundary, and Hamiltonian errors from the executable implementation.
\end{tablenotes}
\end{threeparttable}
\end{table}

The benchmark is passed only when an independently evaluated policy reproduces these targets under the same horizon and normalization. A low HJB residual alone is insufficient; the audit computes the value error, policy error, terminal/boundary error, and maximized-Hamiltonian gap separately. The original figure is retained below as a historical illustration and is explicitly excluded from the R2 evidence ledger.

\subsection{Essential Test: Constrained and Nonsmooth Benchmarks}
The no-short-sale parameterization with $\mu=.01$, $r=.02$, $\gamma=2$, and $\sigma=.2$ has unconstrained share $-0.125$ and constrained optimum $\pi=0$ at every positive wealth level. It is a useful feasibility test, but it has no wealth-dependent switching point and therefore is not a free-boundary or kink benchmark. The revised paper does not call it one.

The nonsmooth diagnostic is instead the exit-time problem in Example \ref{ex:r2_viscosity}. The implementation evaluates the exact supremized operator on a fixed boundary grid and reports the value error at the kink, the boundary error, and the residual profile. A smooth network's localized curvature is an approximation choice; it is not evidence of viscosity selection unless the independently solved benchmark and the boundary audit agree.

\subsection{Benchmarking against Grid-Based Methods: The Scalability Frontier}
The original $d$-stock experiment is separable under its stated primitives: $V(k_1,\ldots,k_d)=\sum_i v_i(k_i)$. It cannot identify an intrinsic high-dimensional advantage. R2 retains it as a separability control and adds a genuinely coupled test with a common resource constraint, correlated shocks, and a cross-stock production term. Both NBO and the sparse-grid baseline receive the same domain, tolerance, hardware class, stopping rule, and held-out test distribution.

The estimands are fixed-accuracy wall-clock time, peak memory, failure rate, value error, policy error, and unilateral Hamiltonian gap. Derivative cost, network size, sample count, optimizer iterations, and probe count are recorded separately. No polynomial complexity claim is made from per-iteration differentiation; the observed scaling is reported with uncertainty across seeds. The methodological contribution is the actor-factorized constrained operator iteration and its auditable decomposition, not mesh-free approximation by itself. Ablations remove actor factorization, hard feasibility, recursive value dependence, and exact-versus-stochastic trace handling.

\subsection{Analysis of Robustness and Reproducibility}

A critical property of a reliable numerical method is its robustness to initialization, sampling, step size, and residual conditioning. R2 reports all declared seeds, including failed pilots, together with stopping criteria and run identifiers. Cross-seed agreement is descriptive evidence for the audited benchmark; it is not evidence of global optimality or of a theorem for Adam. The primary outcome is a vector of errors (value, policy, boundary, and improvement gap), not the HJB residual alone.

\section{Application 1: Recursive Utility and the Structural Stability of NBO}
\label{sec:recursive_test}

A primary motivation for the NBO architecture is to handle recursive utility, where the policy-wise value process is an argument of the flow aggregator. This structure requires a declared recursive utility equation and a value-domain condition; it is not reduced to a generic moving-target slogan.

In this section, we apply NBO to a canonical portfolio choice problem with continuous-time EZW preferences. We provide an analytical target calculation and a protocol for matched baseline comparisons. The historical stability figure is retained as a negative control; the R2 evidence claim rests on the corrected aggregator, domain checks, and run-level provenance.

\subsection{Setup: Epstein-Zin Preferences and the Recursive HJB Equation}
We use the difference-form Epstein--Zin aggregator on the declared positive domain $c>0$ and $bV>0$, with $a=1-1/\psi$ and $b=1-\gamma$:
\begin{equation}
 f_{EZ}(c,V)=\frac{\rho}{a}\left[c^a(bV)^{1-a/b}-bV\right],
\label{eq:ez_aggregator}
\end{equation}
with the logarithmic limit used when $a=0$ and the parameter restrictions needed for the fractional powers stated in the implementation manifest. Its consumption derivative is positive on the declared domain. The value architecture enforces $bV>0$ by a monotone transform, and held-out domain checks verify that no evaluation leaves this domain.

For the stationary homothetic diagnostic, the candidate ratios are
\[
\pi^*=\frac{\mu-r}{\gamma\sigma^2}=0.3,
\qquad m^*=\psi\rho+(1-\psi)\left[r+\frac{(\mu-r)^2}{2\gamma\sigma^2}\right]=0.0455
\]
for the audit tuple $(\rho,\gamma,\psi,\mu,r,\sigma)=(.04,5,1.5,.08,.02,.2)$. These are formal candidates until admissibility, integrability, and verification are checked. They replace the earlier sign-inconsistent formula and the reported value $0.0278$.

\subsubsection{Derivation of the Recursive HJB Equation}
To apply the NBO method, we derive the HJB equation that the value function $V(t,X_t) \equiv J_t$ must satisfy. By applying It\^o's Lemma and the principle of optimality, the drift of the value process plus the aggregator must equal zero for the optimal policy. This yields the recursive HJB equation:
\begin{equation} \label{eq:ez_hjb_full}
-\partial_t V = \sup_{c,\pi}\left\{f_{EZ}(c,V)+(rX+\pi(\mu-r)X-c)V_X+\frac12\pi^2\sigma^2X^2V_{XX}\right\},
\end{equation}
where $\mathcal{D}^{c,\pi}V = (rX + \pi(\mu-r)X - c)\partial_X V + \frac{1}{2}\pi^2\sigma^2X^2 \partial_{XX}^2 V$.

Crucially, the value function $V$ appears directly in the maximand. Furthermore, the optimal portfolio $\pi^*$ depends on the ratio of the Hessian to the gradient, $\pi^* \propto -\frac{V_X}{X V_{XX}}$. This dependence on the second derivative makes the problem particularly sensitive to numerical instability in regimes where $V_{XX}$ is small or noisy.

\subsection{Benchmarking the Recursive Operator}
The recursive experiment is evaluated against the corrected aggregator in Eq. \eqref{eq:ez_aggregator}. We report domain violations, value and policy errors relative to the homothetic diagnostic, and the improvement gap. The historical stability illustration is retained as an archival figure; its constructed comparison is not used as evidence for a competing method. A matched baseline, identical seeds, and deterministic data-generation script are required for any claim about relative stability.

\section{Application 2: Endogenous Preferences and Portfolio Choice}
\label{sec:application_ndu}

We now apply the NBO method to solve a model that is central to the motivating research of \cite{qi2025} but intractable with standard methods: the Merton portfolio problem augmented with Non-Dichotomous Utility (NDU), where an agent can actively manage their own risk aversion. This application serves two purposes. First, it tests the operator interface on a declared two-state preference-formation model. Second, it generates conditional comparative-static hypotheses about state-dependent policies; it is not by itself an identification result or evidence of generic high-dimensional performance.

\subsection{Formal Model Setup and HJB Formulation}
The NDU state space is restricted to $u\in[u_{\min},u_{\max}]$ with a reflected or viability-preserving diffusion specified at the boundary, and to $X\in[X_{\min},X_{\max}]$ with an explicit borrowing/solvency rule. The additive Brownian specification is not used outside this declared domain. Utility is cardinally normalized before $u$ becomes a controlled state; consumption units are fixed and the normalization is reported because shifts in utility are no longer innocuous when preferences are endogenous. The two-dimensional model is presented as an economic application and a low-dimensional audit, not as evidence of generic high-dimensional identification.


The agent's state is the two-dimensional vector $\statevec_t = (u_t, X_t) \in \R^+ \times \R^+$, where $u_t$ is the coefficient of relative risk aversion and $X_t$ is financial wealth. The agent has two external controls, consumption $c_t$ and portfolio share in a risky asset $\pi_t$, and one internal control, the preference adjustment effort $\theta_t$.

The state variables evolve according to the following system of Stochastic Differential Equations (SDEs):
\begin{align}
du_t &= \theta_t dt + \sigma_u dZ_t^u, \label{eq:ndu_u_sde}\\
dX_t &= \left[ (\pi_t (\mu-r) + r)X_t - c_t \right]dt + \pi_t X_t \sigma_X dZ_t^X, \label{eq:ndu_x_sde}
\end{align}
where $Z_t^u$ and $Z_t^X$ are standard Brownian motions with correlation $dZ_t^u dZ_t^X = \rho_{uX} dt$. We assume $\rho_{uX} < 0$, capturing the intuition that adverse financial shocks may directly induce anxiety and a desire for higher risk aversion.

The agent chooses admissible control paths $\{c_s, \pi_s, \theta_s\}_{s=t}^T$ to maximize a time-separable utility functional. The flow utility is of the CRRA-form, but with the risk-aversion coefficient $u_t$ as a state variable. The control of preferences incurs a quadratic cost, $\frac{1}{2}k\theta_t^2$, where $k>0$ is a cost parameter. The agent's objective is:
\begin{equation}
\sup_{\{c_s, \pi_s, \theta_s\}} \E_t \left[ \int_t^T \left( e^{-\rho (s-t)}\frac{c_s^{1-u_s}}{1-u_s} - \frac{1}{2}k\theta_s^2 \right) ds + G(u_T,X_T) \right].
\end{equation}
For clarity in this initial application, we use a time-separable aggregator. The framework's success in Section \ref{sec:recursive_test} demonstrates its capacity to solve more complex NDU models with fully recursive preferences. The choice of a linear control in the drift of the preference state and a quadratic cost of adjustment is standard in the literature on costly adjustment. It provides the simplest canonical representation of a trade-off where changing a state variable is possible but costly, with marginal costs increasing in the size of the adjustment.

Let $V(t, u, X)$ be the value function. Assuming sufficient smoothness, it must satisfy the associated Hamilton-Jacobi-Bellman (HJB) equation. Let $V_t, V_u, V_X, V_{uu}, V_{XX}, V_{uX}$ denote the partial derivatives of $V$. The HJB equation is:
\begin{equation} \label{eq:ndu_hjb}
-\partial_t V = \sup_{c, \pi, \theta} \left\{ \frac{c^{1-u}}{1-u} - \frac{1}{2}k\theta^2 + \mathcal{D}V - \rho V \right\},
\end{equation}
where $\mathcal{D}V$ is the Itô differential operator applied to $V$:
\begin{align}
\begin{split}
\mathcal{D}V = V_u \theta &+ V_X \left[ (\pi(\mu-r) + r)X - c \right] \\
&+ \frac{1}{2}V_{uu}\sigma_u^2 + \frac{1}{2}V_{XX}\pi^2 X^2 \sigma_X^2 + V_{uX} \pi X \sigma_X \sigma_u \rho_{uX}.
\end{split}
\end{align}
The term $\mathcal{H} = \{ \frac{c^{1-u}}{1-u} - \frac{1}{2}k\theta^2 + \mathcal{D}V - \rho V \}$ is the Hamiltonian under the time-relative discount convention. The NBO algorithm seeks policy and value functions that satisfy this equation together with the declared boundary and admissibility conditions.

\subsection{Analytical Derivations of Optimal Policies and Economic Intuition}
While a closed-form solution for $V^*$ is unavailable, we can characterize the optimal policies in terms of the value function and its derivatives by solving the static maximization problem on the right-hand side of the HJB equation \eqref{eq:ndu_hjb}. This provides the rigorous microfoundation for the economic mechanisms we observe in the numerical solution. The first-order conditions (FOCs) with respect to the controls $(c, \theta, \pi)$ are:
\begin{align}
\frac{\partial \mathcal{H}}{\partial c} &= c^{-u} - V_X = 0 \\
\frac{\partial \mathcal{H}}{\partial \theta} &= -k\theta + V_u = 0 \\
\frac{\partial \mathcal{H}}{\partial \pi} &= V_X(\mu-r)X + V_{XX}\pi X^2 \sigma_X^2 + V_{uX} X \sigma_X \sigma_u \rho_{uX} = 0
\end{align}
Solving these FOCs yields the optimal policies as functions of the state and the (unknown) value function derivatives:
\begin{align}
c^*(u, X) &= V_X^{-1/u} \label{eq:ndu_foc_c} \\
\theta^*(u, X) &= \frac{1}{k} V_u \label{eq:ndu_foc_theta} \\
\pi^*(u, X) &= -\frac{V_X(\mu-r)}{V_{XX}X\sigma_X^2} - \frac{V_{uX}\sigma_u \rho_{uX}}{V_{XX}X\sigma_X} \label{eq:ndu_foc_pi}
\end{align}

These equations provide sharp economic predictions:
1.  \textbf{Consumption (Eq. \ref{eq:ndu_foc_c}):} The optimal consumption rule is a generalization of the standard Euler equation, where consumption is a decreasing function of the shadow price of wealth, $V_X$.
2.  \textbf{Preference Adjustment (Eq. \ref{eq:ndu_foc_theta}):} This is the key internal margin. The agent adjusts risk aversion in proportion to the shadow value $V_u$ subject to the declared bounds. For the flow term $c^{1-u}/(1-u)$, $\partial_u[c^{1-u}/(1-u)]=c^{1-u}[1-(1-u)\log c]/(1-u)^2$, which is positive at $(c,u)=(2,2)$. The sign of the complete $V_u$ is therefore an equilibrium object, not an implication of the primitive flow term. Comparative statics in $k$ are reported only after re-solving the model; $\theta=V_u/k$ alone is not a sign theorem.
3.  \textbf{Portfolio Choice (Eq. \ref{eq:ndu_foc_pi}):} The optimal portfolio rule consists of two terms. The first is the classic Mertonian mean-variance demand. The second is a novel \textit{intertemporal hedging demand} against fluctuations in risk aversion. Since we assume $\rho_{uX} < 0$, and expect $V_{uX} > 0$ (the marginal value of wealth is higher when you are more risk-averse), this hedging term is negative, indicating that the agent reduces their risky asset holdings to hedge against shocks that simultaneously reduce wealth and increase risk aversion.

\subsection{Numerical Results and Economic Insights}
The historical policy surfaces are retained as qualitative illustrations. An authoritative R2 claim requires the executable two-dimensional collocation comparison and the error vector in the replication manifest; the figures alone are not treated as evidence that the HJB system has been solved.

The historical Figure \ref{fig:ndu_path} illustrates the proposed mechanism; it is retained as a negative control because its embedded path arrays do not carry R2 run provenance. When wealth is low (left panel), the agent acts to preserve capital. The marginal utility of wealth, $V_X$, is high, leading to low consumption. The agent maintains a high level of risk aversion ($u_t$), which, via Eq. \ref{eq:ndu_foc_pi}, results in a conservative portfolio. In contrast, an agent starting with high wealth (right panel) has a lower marginal utility of wealth. This increases the relative benefit of higher returns, raising the shadow value of having lower risk aversion ($V_u$ becomes more negative). As predicted by Eq. \ref{eq:ndu_foc_theta}, the agent chooses $\theta_t < 0$ to actively drive down $u_t$. This endogenously shaped preference for lower risk aversion then rationalizes the significant increase in the optimal portfolio share $\pi_t$.

\begin{figure}[H]
\centering
\begin{tikzpicture}
% [Code for Figure 3: ndu_path_low and ndu_path_high]
% This code is identical to the original submission.
\begin{groupplot}[
    group style={
        group size=2 by 1,
        xlabels at=edge bottom,
        ylabels at=edge left,
        horizontal sep=2cm,
        vertical sep=1.5cm,
        group name=mygroup
    },
    height=6cm, width=7cm,
    xlabel={Time (Years)},
    xmin=0, xmax=10,
    legend style={at={(0.5,-0.3)}, anchor=north, legend columns=-1}
    ]

    \nextgroupplot[title={Low Initial Wealth ($X_0=0.5$)}, ymin=0]
        \addplot[blue, smooth] table[x=time, y=wealth] {ndu_path_low.dat}; \addlegendentry{Wealth ($X_t$)}
        \addplot[red, dashed, smooth] table[x=time, y=u] {ndu_path_low.dat}; \addlegendentry{Risk Aversion ($u_t$)}
        \addplot[green!50!black, dotted, smooth] table[x=time, y=pi] {ndu_path_low.dat}; \addlegendentry{Portfolio Share ($\pi_t$)}

    \nextgroupplot[title={High Initial Wealth ($X_0=5.0$)}, ymin=0]
        \addplot[blue, smooth] table[x=time, y=wealth] {ndu_path_high.dat}; \addlegendentry{Wealth ($X_t$)}
        \addplot[red, dashed, smooth] table[x=time, y=u] {ndu_path_high.dat}; \addlegendentry{Risk Aversion ($u_t$)}
        \addplot[green!50!black, dotted, smooth] table[x=time, y=pi] {ndu_path_high.dat}; \addlegendentry{Portfolio Share ($\pi_t$)}
\end{groupplot}
\end{tikzpicture}
\caption{\textbf{Simulated Paths of State and Policy Variables.} The figure shows two historical simulated sample paths, visualized by evaluating the trained policy networks over time. The agent actively reduces their risk aversion ($u_t$, red dashed line) when wealth ($X_t$, blue solid line) is high, enabling a much larger portfolio share in the risky asset ($\pi_t$, green dotted line). This is consistent with the derived first-order conditions. This historical illustration is retained for provenance and is not an R2 experimental observation; authoritative R2 outputs are identified by run ID and checksum in the replication manifest. The historical timing is retained for provenance; R2 timings are measured by the declared harness and hardware record.}
\label{fig:ndu_path}
\end{figure}

Figure \ref{fig:ndu_policy} visualizes the learned global policy function for the portfolio share, $\pi^*(u, X)$, by evaluating the trained network on a fine grid of the state space. It reveals a strong "house money" effect mediated by endogenous preferences: as wealth increases (moving right along the x-axis), the agent finds it optimal to take on significantly more risk (the color shifts from blue to yellow). This is not simply due to changing wealth, but because higher wealth incentivizes the agent to actively lower their risk aversion $u_t$, which in turn makes a higher portfolio share $\pi_t$ optimal. This visualization confirms the state-dependent nature of the policies derived in Eqs. \ref{eq:ndu_foc_c}-\ref{eq:ndu_foc_pi}.

\begin{figure}[H]
\centering
\begin{tikzpicture}
% [Code for Figure 4: ndu_policy_surface]
% This code is identical to the original submission.
\begin{axis}[
    title={Learned Optimal Portfolio Policy $\pi^*(u, X)$ Reveals Strong Wealth Effects},
    xlabel={Wealth ($X_t$)},
    ylabel={Risk Aversion ($u_t$)},
    view={0}{90},
    width=12cm,
    height=9cm,
    colorbar,
    colorbar style={
        ylabel={Portfolio Share $\pi^*(u,X)$}
    },
    colormap/viridis,
    ]
    \addplot3[surf, shader=faceted] file {ndu_policy_surface.dat};
\end{axis}
\end{tikzpicture}
\caption{\textbf{Learned Optimal Policy in the NDU Portfolio Model.} The figure shows a surface plot of the learned optimal portfolio share $\pi^*(u_t, X_t)$, visualized by evaluating the trained policy network on a fine grid of the state space. The color indicates the optimal allocation to the risky asset. The strong positive relationship between wealth and risk-taking is mediated by the agent's endogenous choice of their own risk aversion. This historical illustration is retained for provenance and is not an R2 experimental observation; authoritative R2 outputs are identified by run ID and checksum in the replication manifest.}
\label{fig:ndu_policy}
\end{figure}


\subsection{Comparative Statics and Implications for Structural Estimation}
To further demonstrate the model's economic content and the solver's utility, we conduct a comparative statics exercise on the preference adjustment cost parameter, $k$. As can be seen from the first-order condition, $\theta^* = (1/k)V_u$, the adjustment cost $k$ directly dampens the agent's response to any given incentive to change their preferences.

The historical Figure \ref{fig:ndu_statics} provides a qualitative illustration of this analytical result. It shows the optimal portfolio share as a function of wealth for a fixed level of risk aversion ($u=2$), but under two different cost regimes: low cost ($k=1.0$) and high cost ($k=5.0$). The results are intuitive and economically meaningful. When the cost of adjustment is low (blue line), the agent aggressively increases portfolio risk as wealth rises, reflecting their ability to cheaply lower their risk aversion to justify the higher exposure. When the cost is high (red line), this effect is significantly muted. This is a comparative-static illustration conditional on the declared state constraints and independent solver check; it does not establish identification of $k$.

\begin{figure}[H]
\centering
\begin{tikzpicture}
% [Code for Figure 5: ndu_statics]
% This code is identical to the original submission.
\begin{axis}[
    title={\textbf{Comparative Statics: Effect of Adjustment Cost $k$}},
    xlabel={Wealth ($X_t$)},
    ylabel={Optimal Portfolio Share $\pi^*(X_t | u_t=2)$},
    width=10cm,
    height=7cm,
    legend pos=north west,
    ]

    \addplot[thick, blue, smooth, mark=*, mark options={solid}] table[x=wealth, y=pi_low_k] {ndu_statics.dat};
    \addlegendentry{Low Cost ($k = 1.0$)}

    \addplot[thick, red, dashed, smooth, mark=square*, mark options={solid}] table[x=wealth, y=pi_high_k] {ndu_statics.dat};
    \addlegendentry{High Cost ($k = 5.0$)}

\end{axis}
\end{tikzpicture}
\caption{\textbf{Effect of Preference Adjustment Cost.} The figure shows the optimal portfolio share as a function of wealth, conditional on the agent's current risk aversion being $u_t=2$. A higher cost of adjusting preferences ($k$) significantly dampens the agent's willingness to increase risk-taking as wealth grows, confirming the mechanism in Eq. \ref{eq:ndu_foc_theta}. This historical illustration is retained for provenance and is not an R2 experimental observation; authoritative R2 outputs are identified by run ID and checksum in the replication manifest.}
\label{fig:ndu_statics}
\end{figure}

The dependence on $k$ motivates a possible structural-estimation exercise, but identification requires an observation model, a specified moment map, and a rank or injectivity argument. The present figures supply a simulation-based comparative-static illustration only. A future structural-estimation exercise could embed the audited solver in simulated method of moments, but that exercise requires an observation model, an identified moment map, and an independent treatment of latent preference states. Those ingredients are outside the present numerical illustration.

\section{Application 3: Time-Inconsistency and Dynamic Games}
\label{sec:application_extensions}

A significant strength of the NBO method is its flexibility. The core actor-critic architecture, stabilized by the HJB loss, can be adapted to solve a much broader class of problems beyond single-agent, time-consistent control. This section moves beyond the theoretical frameworks outlined in the appendix to provide concrete numerical applications, demonstrating the practical utility of the NBO method for finding equilibria in models with time-inconsistent preferences and in high-dimensional dynamic games.

\subsection{Sophisticated Quasi-Hyperbolic Agents in a Consumption-Savings Problem}
We define a temporal self by the one-shot deviation objective
\begin{equation}
\sup_c\{u(c)+\beta\mathcal V^{\pi}(s')\},
\label{eq:r2_soph_objective}
\end{equation}
where the continuation policy is re-optimized by future selves. The equilibrium policy is therefore characterized by an extended HJB system, not by a target-policy lag alone. The implementation checks the one-shot deviation gain on a held-out state grid and verifies the no-present-bias limit $\beta=1$ against the exponential-discount solution. Since $\beta<1$ lowers the current utility weight in the local FOC $\beta u'(c)=V_X$, a claim of over-consumption is made only after the induced change in $V_X$ is measured; it is not inferred while holding the shadow price fixed.

The main text and supplement now use this one authoritative equilibrium system. Consumption columns are explicitly labeled as levels or consumption--wealth ratios, and the continuation dynamics include the discount kernel used in the objective. Historical target-policy and local-Hamiltonian formulations remain available in the archive but are not combined as if they were one model.

\subsection{Markov Perfect Equilibrium in a Dynamic Cournot Game}
For player $i$, the improvement step differentiates player $i$'s Hamiltonian with respect to player $i$'s parameters while detaching $\pi_{-i}$. The equilibrium audit reports the unilateral exploitability gap
\[
\mathsf{Exploit}_i=\sup_{\tilde\pi_i}\{J_i(\tilde\pi_i,\pi_{-i})-J_i(\pi_i,\pi_{-i})\},
\]
computed by an independent best-response solver. A joint optimum is not a Nash equilibrium: in the static benchmark $u_i=q_i(1-q_i-q_j)$, symmetric Nash quantities are $1/3$, whereas maximizing joint profit gives $1/4$ and leaves a profitable unilateral deviation. This benchmark is a regression test for the derivative graph.

Comparisons across firm counts use a market-size normalization that keeps aggregate demand and per-firm capital comparable. The $N\to\infty$ limit is reported as a separately solved competitive benchmark, with investment and depreciation costs retained. Symmetry of initialization is not treated as an equilibrium-selection argument; results are reported across multiple initializations and with exploitability checks.

\section{Conclusion} \label{sec:conclusion}
This revision develops NBO as a separated, auditable approximation to controlled Bellman operators. The central contribution is the explicit decomposition into policy evaluation, feasible policy improvement, boundary verification, and equilibrium diagnostics. That decomposition preserves the method's ability to accommodate recursive utility, endogenous preferences, temporal selves, and strategic interaction while removing the invalid identification of a joint penalty minimizer with the economic solution.

The mathematical result is conditional on a declared Markov boundary-value problem, comparison, exact evaluation, and pointwise improvement. The finite-step neural implementation is assessed by held-out value, policy, boundary, and exploitability errors. The corrected Merton and Epstein--Zin anchors, the exit-time nonsmooth diagnostic, and the player-specific Cournot test are part of the regression suite. The endogenous-preference and time-inconsistency applications remain substantive economic exercises, but their comparative statics and identification content are stated only when the model and observation map support them. High-dimensional claims are tied to coupled, separability-aware benchmarks and matched resource budgets.

\paragraph*{Replication Note.} All authoritative R2 sources, configurations, raw-output schemas, checksums, and unsuccessful pilots are listed in the revision manifest. Historical TeX sources and illustrative figures are retained unchanged for provenance and are labeled as such; they do not supply R2 numerical evidence. The source package can be rebuilt with the command in the replication README after installing the declared Econometric Society and LaTeX dependencies.

\section{Numerical and Implementation Details}
\label{app:implementation}

This appendix provides a formal account of the numerical methods and implementation choices underpinning the NBO algorithm, ensuring full transparency and reproducibility.

\subsection{Network Architecture and Optimization}
The NBO framework utilizes three feedforward neural networks (multi-layer perceptrons) to approximate the value and policy functions. Let the full set of trainable parameters be $\paramvec = (\psi, \phi, \xi)$.
\begin{itemize}
    \item \textbf{Value Network ($V_\xi$):} Maps the state and time $(t, \statevec_t) \in [0,T] \times \mathbb{R}^d$ to the scalar value function estimate, $V_\xi(t, \statevec_t) \in \mathbb{R}$.
    \item \textbf{Internal Policy Network ($\Theta_\psi$):} Maps $(t, \statevec_t)$ to the internal control vector, $\theta_t = \Theta_\psi(t, \statevec_t) \in \mathbb{R}^q$.
    \item \textbf{External Policy Network ($\mathcal{A}_\phi$):} Maps $(t, \statevec_t)$ to the external control vector, $a_t = \mathcal{A}_\phi(t, \statevec_t) \in \mathbb{R}^k$.
\end{itemize}
Each network consists of $L$ hidden layers, each with $N_h$ neurons. For all experiments in this paper, we use $L=3$ layers with $N_h=256$ neurons. For $l \in \{1, \dots, L\}$, the transformation at layer $l$ is $\bm{h}_l = g(\bm{W}_l \bm{h}_{l-1} + \bm{b}_l)$, where $\bm{h}_{l-1}$ is the output of the previous layer, $(\bm{W}_l, \bm{b}_l)$ are the weights and biases of layer $l$, and $g(\cdot)$ is a non-linear activation function. To ensure theoretical consistency with the requirements of our HJB-based loss function (see the derivative-approximation discussion in the theory appendix), we use the Gaussian Error Linear Unit (GELU), $g(x) = x \Phi(x)$ where $\Phi$ is the standard normal CDF, for all hidden layers. The GELU activation function is chosen not primarily for performance but to satisfy the theoretical conditions of the derivative-approximation result, which underpins the validity of our entire HJB-based optimization framework, as it is infinitely differentiable ($C^\infty$) and thus satisfies the necessary smoothness assumptions. The full set of weights and biases constitutes the parameter vectors $\psi, \phi, \xi$.

The finite-sample implementation uses Adam separately for the critic evaluation parameters and the actor improvement parameters. The critic step differentiates $L_E$ with respect to $\xi$ while the actor step differentiates $L_I$ with respect to $(\psi,\phi)$ with $V_\xi$ detached. For player $i$, rival policy outputs are detached. Adam's constant-step updates are recorded as finite-sample evidence and are not covered by the conditional stochastic-approximation statement.

\subsection{SDE Discretization and Monte Carlo Estimation}
The expectations in the loss functions are approximated via Monte Carlo integration over trajectories of the state variables. These trajectories are generated by discretizing the governing SDEs. Let the time horizon $[0, T]$ be divided into $N$ steps of size $\Delta t = T/N$, creating a time grid $t_k = k\Delta t$ for $k=0, \dots, N$. The continuous-time state process is
\begin{equation}
d\statevec_t = \driftvec_{\statevec}(t, \statevec_t, \controlvec_t) dt + \diffvec_{\statevec}(t, \statevec_t, \controlvec_t) d\bm{M}_t,
\end{equation}
where $\controlvec_t = (\Theta_\psi(t,\statevec_t), \mathcal{A}_\phi(t,\statevec_t))$ is the policy function, and $\bm{M}_t$ is a continuous vector-valued martingale with quadratic variation process $d\langle \bm{M} \rangle_t = \bm{Q}(t, \statevec_t) dt$.

We use the Euler-Maruyama scheme for discretization. The state at time $t_{k+1}$ is approximated as:
\begin{equation}
\statevec_{t_{k+1}} \approx \statevec_{t_k} + \driftvec_{\statevec}(t_k, \statevec_{t_k}, \controlvec_{t_k}) \Delta t + \diffvec_{\statevec}(t_k, \statevec_{t_k}, \controlvec_{t_k}) \Delta \bm{M}_{t_k}.
\end{equation}
The martingale increment $\Delta \bm{M}_{t_k} = \bm{M}_{t_{k+1}} - \bm{M}_{t_k}$ is a random variable sampled from a distribution with mean zero and covariance given by the discretized quadratic variation:
\begin{equation}
\E[\Delta \bm{M}_{t_k} \Delta \bm{M}_{t_k}^\top \mid \mathcal{F}_{t_k}] = \E \left[ \int_{t_k}^{t_{k+1}} \bm{Q}(s, \statevec_s) ds \mid \mathcal{F}_{t_k} \right] \approx \bm{Q}(t_k, \statevec_{t_k}) \Delta t.
\end{equation}
For the declared Brownian Markov specification, we sample $\Delta \bm{M}_{t_k}$ from $\mathcal{N}(\bm{0},\bm{I}\Delta t)$ and apply the diffusion loading. A general continuous martingale is not replaced by Gaussian increments merely from its quadratic variation; such a model requires a separate approximation argument. The integrals in the loss functions are then approximated by sums over the discretized trajectory, and the expectation is approximated by the sample mean over a batch of $M$ such simulated trajectories (we use $M=1024$).

\subsection{Computation of Derivatives via Automatic Differentiation}
A critical component of the NBO method is the computation of the value function's derivatives, $\partial_t V_\xi, \nabla_{\statevec} V_\xi, \nabla_{\statevec\statevec}^2 V_\xi$, which are required for the HJB loss. These are computed using automatic differentiation (AD), a set of techniques for numerically evaluating the derivative of a function specified by a computer program. AD leverages the fact that any program executes a sequence of elementary arithmetic operations. By applying the chain rule repeatedly to these operations, one can compute derivatives automatically and to machine precision.

The gradients $\partial_t V_\xi$ and $\nabla_{\statevec} V_\xi$ are computed efficiently via a single pass of the backpropagation algorithm. The Hessian matrix $\nabla_{\statevec\statevec}^2 V_\xi$ is more computationally intensive. Modern deep learning libraries (like PyTorch or TensorFlow) compute the term $\text{Tr}(\bm{\sigma} \bm{Q} \bm{\sigma}^\top \nabla_{\statevec\statevec}^2 V)$ efficiently, often without explicitly constructing the full Hessian matrix. This is typically achieved using Hessian-vector products, which scales more favorably with the state dimension than forming the full Hessian.

\begin{appendix}
\section{Hyperparameter Selection and Sensitivity}
The separated implementation reports the critic boundary weight $\lambda_B$, actor constraint weight $\lambda_g$, actor and critic step sizes, batch size, sampling measure, and Hessian mode. Sensitivity is evaluated by held-out value, policy, boundary, residual-quantile, and improvement-gap metrics. These quantities are tuning parameters for a finite-sample approximation; they do not create the step-size separation required by the asymptotic theorem.

The old $\omega$ sensitivity curves are retained only as a negative control for the historical composite objective. They are not used to support a claim about the R2 algorithm. The authoritative sensitivity record, when a training run is supplied, is a JSONL table conforming to \texttt{replication/results\_schema.json} with failed pilots retained.

\section{Formal Proofs}
\label{app:proofs}
This appendix proves the conditional statement actually used in R2. Under Assumption \ref{ass:r2_markov}, exact policy evaluation solves \eqref{eq:r2_policy_eval}. Pointwise policy improvement gives
\[
\mathcal H^{\pi_{n+1}}[V^{\pi_n}](t,s)\geq \mathcal H^{\pi_n}[V^{\pi_n}](t,s)
\]
for every $(t,s)$. The comparison principle and the standard verification inequality then imply monotonicity of the evaluated value sequence whenever the policy class is closed under the improvement map. Compactness gives accumulation points; continuity of the operator passes the improvement inequality to any accumulation point. At an accumulation point, the evaluation equation and the pointwise improvement condition jointly give the HJB equation. Uniqueness follows from comparison.

No step in this proof applies to the historical joint loss, to a finite network with nonzero residual, or to constant-step Adam. The two counterexamples in Section \ref{sec:theoretical_framework} are included to make those exclusions testable.

\section{Theoretical Underpinnings of the NBO Algorithm}
\label{app:uat_conv}
Smooth neural networks approximate $C^2$ value functions and their derivatives on compact sets under the usual nonpolynomial-activation conditions. This fact is used only on regions where the target is classically twice differentiable. It does not approximate a derivative at a kink and does not select a viscosity solution from an $L^2$ residual.

\subsection{Conditional stochastic-approximation statement}
For projected updates with step sizes $a_n,b_n$ satisfying $\sum a_n=\sum b_n=\infty$, $\sum(a_n^2+b_n^2)<\infty$, $a_n/b_n\to0$, bounded iterates, martingale-difference noise with uniformly bounded second moments, and a globally attracting critic ODE for each frozen policy, the two-timescale ODE method yields convergence to the internally chain-transitive set of the actor ODE. If that set consists of isolated stationary points and the improvement gap is zero only at optimal policies, the limit is a policy-iteration fixed point. These are hypotheses for a separate asymptotic algorithm; the Adam runs in this paper use constant step sizes and are therefore finite-sample evidence only.

\section{Robustness and Reproducibility}
The R2 audit records initialization, seed, sampling measure, step sizes, boundary weight, constraint parameterization, Hessian mode, stopping rule, runtime, peak memory, and failure status for every run. It reports held-out value error, policy error, boundary error, residual quantiles, and unilateral improvement or exploitability gaps. No run is omitted because it failed, and no constructed curve is presented as an experimental trajectory.

\end{appendix}



\phantomsection
\addcontentsline{toc}{section}{References}
\bibliography{reference}
\bibliographystyle{ecta-fullname}


\clearpage
\section*{Supplementary Appendix to ``Neural Bellman Operators'' (Not For Publication)}\label{sec:olappendix}
\setcounter{equation}{0}
\renewcommand{\theequation}{S\arabic{equation}}
\setcounter{section}{0}
\renewcommand{\thesection}{S\arabic{section}}

\section{Time-Inconsistent Agents: Extended-HJB Specification}
A sophisticated temporal self chooses a one-shot action while taking the continuation policy of future selves as given. Let $\mathcal V$ denote the continuation value under that policy. The equilibrium condition is
\begin{equation}
 c^*(t,s)\in\argmax_{c\in\mathcal C(s)}\{\beta u(c)+\mathcal D^{c}\mathcal V(t,s)\},
\qquad
-\mathcal V_t=u(c^*)+\mathcal D^{c^*}\mathcal V,
\label{eq:supp_r2_ehjb}
\end{equation}
with terminal and boundary data stated by the application. The policy-evaluation step uses the second equation with $c^*$ detached; the improvement step uses the first equation with $\mathcal V$ detached. The equilibrium residual is the one-shot deviation gain evaluated on a held-out grid, not a penalty that silently combines the two graphs. The no-present-bias limit $\beta=1$ is tested against the exponential-discount solution.

\section{Dynamic Games: Player-Specific Operator Iteration}
For each player $i$, rivals $\pi_{-i}$ are detached during the improvement step:
\begin{equation}
\pi_i^{+}\in\argmax_{a_i\in\mathcal A_i}\{f_i(s,a_i,\pi_{-i})+\mathcal D^{a_i,\pi_{-i}}V_i\}.
\label{eq:supp_r2_game}
\end{equation}
The evaluation equation is solved for $V_i$ at the frozen profile. The audit reports the unilateral gain from an independently computed best response and repeats the calculation from multiple initial profiles. In the static Cournot regression test, $q_i=1/3$ is the Nash quantity while $q_i=1/4$ is the joint-profit quantity; the latter therefore fails the unilateral test.

\section{Evidence and Negative Controls}
All supplement experiments use the schemas in \texttt{replication/results\_schema.json}. Every table row and plot curve carries a run identifier, seed, configuration hash, and source checksum. Historical figures are kept as negative controls when their source arrays are inconsistent with the stated experiment. They are never used to support a new numerical claim.
\end{document}
